% !TEX program = xelatex
\documentclass[10pt,a4paper]{article}

% ============================================================
% COMPILATION
% ============================================================
% Compile with XeLaTeX so that Telugu and Devanagari examples
% can be displayed correctly.

% ============================================================
% PACKAGES
% ============================================================
\usepackage[a4paper,margin=0.62in]{geometry}
\usepackage{graphicx}
\usepackage{booktabs}
\usepackage{multirow}
\usepackage{array}
\usepackage{amsmath}
\usepackage{float}
\usepackage{subcaption}
\usepackage{enumitem}
\usepackage{microtype}
\usepackage{hyperref}
\usepackage{xurl}
\usepackage{titlesec}
\usepackage{setspace}
\usepackage{fontspec}

% ============================================================
% FONTS
% ============================================================
\newfontfamily\telugufont{Noto Sans Telugu}
\newfontfamily\devanagarifont{Noto Sans Devanagari}

% ============================================================
% PAGE / FORMATTING
% ============================================================
\hypersetup{
    colorlinks=true,
    linkcolor=black,
    urlcolor=blue,
    citecolor=black
}

\setstretch{0.95}

\setlength{\parindent}{0pt}
\setlength{\parskip}{3pt}

\setlist[itemize]{
    nosep,
    leftmargin=*
}

\setlist[enumerate]{
    nosep,
    leftmargin=*
}

\titlespacing*{\section}{0pt}{6pt}{3pt}
\titlespacing*{\subsection}{0pt}{5pt}{2pt}

\titleformat{\section}
    {\large\bfseries}
    {\thesection}
    {0.5em}
    {}

\titleformat{\subsection}
    {\normalsize\bfseries}
    {\thesubsection}
    {0.5em}
    {}

% ============================================================
% FIGURE HELPER
% ============================================================
\newcommand{\safeimage}[2][\linewidth]{%
    \IfFileExists{#2}
    {\includegraphics[width=#1]{#2}}
    {\fbox{\parbox[c][1.7cm][c]{0.9\linewidth}
    {\centering Figure not uploaded}}}
}

% ============================================================
% TITLE
% ============================================================
\title{
    \vspace{-1.0cm}
    \textbf{Language Models and Agents: Individual Project Report - Bonus}\\
}

\author{
    Korlapati Keerthana\\
    2023101121
}

\date{}

% ============================================================
% DOCUMENT
% ============================================================
\begin{document}

\maketitle
\vspace{-0.5cm}

% ============================================================
\section{Bonus: Positional-Embedding Ablation}
% ============================================================

The optional positional-embedding ablation was performed for
(Model H) Telugu. The ablated model removes positional embeddings
while retaining the rest of the Transformer architecture.

\begin{table}[H]
\centering
\small
\caption{Telugu positional-embedding ablation.}
\begin{tabular}{lcc}
\toprule
Metric & \textbf{Standard} & \textbf{No position}\\
\midrule
Parameters & 23.04M & 22.81M\\
Test loss & 3.124 & 3.171\\
Perplexity & 22.73 & 23.83\\
BPB & 0.467 & 0.474\\
\bottomrule
\end{tabular}
\end{table}

Removing positional embeddings increases test perplexity from 22.73
to 23.83 and BPB from 0.467 to 0.474. This is consistent with the
role of positional information in a decoder-only Transformer: without
explicit position information, self-attention does not directly
represent token order.

\subsection{Bonus Generation Evaluation}

The full generation evaluation was also run for the no-position
model.

\begin{table}[H]
\centering
\scriptsize
\caption{Telugu no-position generation metrics.}
\begin{tabular}{lrrrrrr}
\toprule
Setting & BLEU & chrF & ROUGE-L & Rep. & D-1 & D-2\\
\midrule
Greedy & 8.66 & 23.16 & .141 & .558 & .288 & .392\\
$T=0.5$ & 7.50 & 23.52 & .133 & .264 & .416 & .634\\
$T=1.0$ & 4.78 & 22.50 & .089 & .016 & .732 & .961\\
$T=1.5$ & .55 & 17.22 & .025 & .000 & .811 & .998\\
\bottomrule
\end{tabular}
\end{table}

The generation results are mixed compared with the standard model.
Some reference-overlap metrics are slightly higher for the ablation,
but this does not change the intrinsic language-modeling result.
Repetition is slightly higher for greedy decoding and Distinct-1 is
lower.

The same temperature trend is still visible: higher temperature
reduces repetition and increases diversity, while very high
temperature reduces reference overlap.

\subsection{Bonus Training Curve}

\begin{figure}[H]
\centering
\begin{subfigure}{0.48\linewidth}
    \centering
    \safeimage{figures/bonus_train_vs_validation_loss.png}
    \caption{Training/validation loss}
\end{subfigure}
\hfill
\begin{subfigure}{0.48\linewidth}
    \centering
    \safeimage{figures/bonus_learning_rate_schedule.png}
    \caption{Learning-rate schedule}
\end{subfigure}
\caption{Training results for the Telugu no-position ablation.}
\end{figure}

\subsection{Bonus Attention Analysis}

The no-position model was also evaluated with the same attention
analysis used for the standard model.

\begin{table}[H]
\centering
\small
\caption{Telugu no-position attention statistics.}
\begin{tabular}{lcc}
\toprule
Metric & \textbf{Standard} & \textbf{No position}\\
\midrule
Mean attention entropy & 2.187 & 2.304\\
Mean attention distance & 11.78 & 11.45\\
\bottomrule
\end{tabular}
\end{table}

\begin{figure}[H]
\centering

\begin{subfigure}{0.235\linewidth}
    \centering
    \safeimage{figures/bonus_attention_early_layer_00_head_00.png}
    \caption{Early H0}
\end{subfigure}
\hfill
\begin{subfigure}{0.235\linewidth}
    \centering
    \safeimage{figures/bonus_attention_early_layer_00_head_01.png}
    \caption{Early H1}
\end{subfigure}
\hfill
\begin{subfigure}{0.235\linewidth}
    \centering
    \safeimage{figures/bonus_attention_early_layer_00_head_02.png}
    \caption{Early H2}
\end{subfigure}
\hfill
\begin{subfigure}{0.235\linewidth}
    \centering
    \safeimage{figures/bonus_attention_early_layer_00_head_03.png}
    \caption{Early H3}
\end{subfigure}

\vspace{2pt}

\begin{subfigure}{0.235\linewidth}
    \centering
    \safeimage{figures/bonus_attention_late_layer_09_head_00.png}
    \caption{Late H0}
\end{subfigure}
\hfill
\begin{subfigure}{0.235\linewidth}
    \centering
    \safeimage{figures/bonus_attention_late_layer_09_head_01.png}
    \caption{Late H1}
\end{subfigure}
\hfill
\begin{subfigure}{0.235\linewidth}
    \centering
    \safeimage{figures/bonus_attention_late_layer_09_head_02.png}
    \caption{Late H2}
\end{subfigure}
\hfill
\begin{subfigure}{0.235\linewidth}
    \centering
    \safeimage{figures/bonus_attention_late_layer_09_head_03.png}
    \caption{Late H3}
\end{subfigure}

\caption{Attention maps for the Telugu no-position ablation.}
\end{figure}

The no-position model has slightly higher average attention entropy
and slightly lower mean attention distance than the standard model.
The heatmaps still show different attention patterns across heads,
although the model no longer receives an explicit positional
embedding.





\end{document}